% Figure: the falling-domino picture of induction. The base case tips
% the first tile; the inductive step is the guarantee that each tile
% tips its successor.
\begin{figure}[htbp]
\centering
\begin{tikzpicture}[
  tile/.style={draw=engblack, thick, fill=engblue!12,
    minimum width=0.32cm, minimum height=1.3cm},
  fall/.style={draw=engblack, thick, fill=engred!18,
    minimum width=0.32cm, minimum height=1.3cm},
  ann/.style={font=\footnotesize, align=center},
]
% A few fallen tiles (rotated), then standing tiles.
\node[fall, rotate=70] at (0.0, 0.45) {};
\node[fall, rotate=55] at (0.75, 0.55) {};
\node[fall, rotate=40] at (1.55, 0.62) {};
\node[tile, rotate=20] at (2.45, 0.66) {};
\node[tile] at (3.4, 0.65) {};
\node[tile] at (4.1, 0.65) {};
\node[tile] at (4.8, 0.65) {};
\node[ann] at (5.7, 0.65) {$\cdots$};

% ground
\draw[draw=engblack, thick] (-0.6,0) -- (6.2,0);

% labels
\node[ann] at (0.0, -0.55) {$P(1)$};
\node[ann] at (3.4, -0.55) {$P(n)$};
\node[ann] at (4.1, -0.55) {$P(n{+}1)$};

\node[ann, anchor=west] at (-0.6, 2.3)
  {base case: $P(1)$ tips the first tile};
\draw[->, thick, draw=engblack] (1.4,2.1) -- (0.4,1.1);

\node[ann, anchor=west] at (3.2, 2.3)
  {inductive step: each tile tips its successor};
\draw[->, thick, draw=engblack] (4.2,2.1) .. controls (4.0,1.4) .. (3.7,1.2);
\end{tikzpicture}
\caption[Induction as falling dominoes]{Induction in the
falling-domino picture. The base case tips the first tile; the
inductive step is the standing guarantee that a fallen $P(n)$ tile
tips its $P(n+1)$ successor. Either half alone proves nothing: a base
case with no step leaves all later tiles standing, and a step with no
base case leaves the whole row standing. The picture is a mnemonic;
the mathematical content is the induction axiom of the natural
numbers.}
\label{fig:vol02:ch01:induction-dominoes}
\end{figure}
